% Part: applied-modal-logics
% Chapter: epistemic-logic
% Section: introduction

\documentclass[../../../include/open-logic-section]{subfiles}

\begin{document}

\olfileid[pt]{aml}{el}{int}

\olsection{Introdução}

Assim como a lógica modal trata de \emph{proposições modais} e das relações de consequência entre elas, a lógica epistêmica trata de \emph{proposições epistêmicas} e das relações de consequência entre elas. Em vez de interpretar os operadores modais como representações de possibilidade e necessidade, os conectivos unários recebem interpretações epistêmicas ou doxásticas, a fim de modelar conhecimento e crença. Podemos, por exemplo, querer expressar afirmações como estas:

\begin{enumerate}
\item Richard sabe que Calgary fica em Alberta.
\item Audrey considera possível que haja um cachorro no sofá.
\item Richard sabe que Audrey sabe que sua aula é às terças-feiras.
\item Todos sabem que um ano tem 12 meses.
\end{enumerate}
A origem da lógica epistêmica contemporânea é frequentemente associada a \emph{Knowledge and Belief}, de Jaako Hintikka, publicado em 1962, quando a semântica de mundos possíveis passava a ser cada vez mais usada na lógica. De fato, as lógicas epistêmicas empregam a maioria das mesmas ferramentas semânticas que outras lógicas modais, mas as interpretam de modo diferente. A principal mudança está no que tomamos a \emph{relação de acessibilidade} como representando. Nas lógicas epistêmicas, ela representa alguma forma de \emph{possibilidade epistêmica}. Veremos que a noção epistêmica modelada afeta as restrições que desejamos impor à relação de acessibilidade. Veremos também o que ocorre com a teoria da correspondência quando ela recebe uma interpretação epistêmica. Observe que os exemplos acima mencionam dois agentes, Richard e Audrey, e a relação entre aquilo que cada um sabe. As lógicas epistêmicas que consideraremos serão lógicas multiagente, nas quais tais coisas podem ser expressas. Em contraste, uma lógica epistêmica de agente único trataria apenas do que um indivíduo sabe ou acredita.

\end{document}
